\scanpage{39}
Soit maintenant (pour $n\geq1$ fixé) $D$ un
diviseur relatif de rang $n$ sur $F$. Soit
\inserted{$M_D\subset\underline{\mathrm{Sym}}^n(\check V)$} le sous-\deleted{\unsure{fibré de rang|faisceau des}}Module
des fonctions polynômiales sur $F$ de degré $\leq n$
\unsure{s'y|se|s'annulant} \unsure{tuant|trouvant|sur $D$}.

\uline{Proposition} $M_D$ est un supplémentaire de
$\underline{\mathrm{Sym}}^{n-1}(\check V)$ dans $\mathrm{Sym}^n(\check V)$, i.e. le composé
\diagram{\begin{tikzcd}[ampersand replacement=\&,column sep=small]
M_D\arrow[r,hook]\arrow[rrr,bend right=14] \& \mathrm{Sym}^{\deleted{\unsure{n-1}}n}(\check V)\arrow[r,dash] \& \mathrm{Sym}^n(\check V)/\mathrm{Sym}^{n-1}(\check V)\arrow[r,"\sim",no head] \& L^{\otimes -n}
\end{tikzcd}}
est un iso. L'application
\[
D\longrightarrow M_D
\]
est une bijection de l'ens des diviseurs
relatifs sur $F$ de degré $n$ avec l'ens
des sous-\deleted{\unsure{M}}Modules $M_D$ de $\mathrm{Sym}^n(V)$
satisfaisant la condition précédente. Cette
bijection est compatible avec les chgts de
base, et définit un isomorphisme
\diagram{\begin{tikzcd}[ampersand replacement=\&,column sep=small,row sep=small]
\underline{\mathrm{Div}}^{+n}_{F/S}\arrow[r,"\sim"] \& \substack{\text{\uline{\unsure{Splittages|Splittings} de}}\\
(0\longrightarrow\mathrm{Sym}^{n-1}(\check V)\longrightarrow\mathrm{Sym}^n(\check V)\longrightarrow L^{\otimes -n}\longrightarrow\unsure{0})}\\
\underset{\text{puissance sym. nième}}{\Sigma^n(F/S)}\arrow[u,"\sim"] \& \substack{\text{\uline{splittages de} } L^{\otimes\unsure{(+n)|n}}\otimes\text{cette suite}\\
\text{exacte, i.e. de la suite}\\ \inserted{\unsure{donnée|(50)}}\\
\deleted{\unsure{0\longrightarrow L\otimes\mathrm{Sym}^{n-1}(V)\longrightarrow\mathrm{Sym}^n(V)\longrightarrow\mathcal O_S\longrightarrow0}}}\arrow[u,no head,"\sim"]
\end{tikzcd}}
% archive sheet number: 38
